% ID: FIS-URT-COM-001
% Titolo: Urto elastico osservato dal centro di massa
% Disciplina: Fisica
% Area: Quantita di moto e urti
% Argomento: Centro di massa
% Sottoargomento: Trasformazione delle velocita nel sistema del centro di massa
% Classe: Terza liceo scientifico
% Difficolta: 2
% Tipo: problema_numerico
% Risultato: prima +2,6 m/s e -2,6 m/s; dopo -2,6 m/s e +2,6 m/s
% Tempo_stimato: 7 min
% Competenze: cambio di sistema di riferimento, interpretazione dell'urto elastico, uso della velocita del centro di massa
% Prerequisiti: centro di massa, urto elastico unidimensionale, velocita relativa
% Tag: fisica, quantita_moto, centro_di_massa, urti_elastici, sistemi_riferimento
% Fonte: rielaborazione_da_traccia_fornita_dal_docente
% Autore: Riccardo Carli
% Licenza: CC-BY-SA-4.0

\begin{esercizio}
Due palle da biliardo hanno la stessa massa. Nel sistema di riferimento del
tavolo, la prima si muove con velocita \(5{,}2\,\mathrm{m/s}\) verso la
seconda, inizialmente ferma. L'urto e frontale ed elastico.

Determina le velocita delle due palle, prima e dopo l'urto, nel sistema di
riferimento solidale con il centro di massa del sistema.
\end{esercizio}

\begin{soluzione}
Indichiamo con \(m\) la massa di ciascuna palla e scegliamo positivo il verso
del moto iniziale della prima palla. Nel sistema del tavolo:
\[
u_1=5{,}2\,\mathrm{m/s},
\qquad
u_2=0.
\]
La velocita del centro di massa e
\[
v_{\mathrm{CM}}
=\frac{m u_1+m u_2}{2m}
=\frac{u_1+u_2}{2}
=2{,}6\,\mathrm{m/s}.
\]

Nel sistema del centro di massa si sottrae \(v_{\mathrm{CM}}\) a ogni
velocita. Prima dell'urto:
\[
u_1^*=u_1-v_{\mathrm{CM}}=2{,}6\,\mathrm{m/s},
\]
\[
u_2^*=u_2-v_{\mathrm{CM}}=-2{,}6\,\mathrm{m/s}.
\]

In un urto elastico frontale tra masse uguali, nel sistema del tavolo le due
palle si scambiano le velocita:
\[
v_1=0,
\qquad
v_2=5{,}2\,\mathrm{m/s}.
\]
La velocita del centro di massa resta invariata. Dopo l'urto:
\[
v_1^*=v_1-v_{\mathrm{CM}}=-2{,}6\,\mathrm{m/s},
\]
\[
v_2^*=v_2-v_{\mathrm{CM}}=2{,}6\,\mathrm{m/s}.
\]
Nel sistema del centro di massa le due velocita hanno quindi sempre lo stesso
modulo e versi opposti; nell'urto elastico ciascuna palla inverte semplicemente
il proprio verso.
\end{soluzione}
